% Part: normal-modal-logic
% Chapter: syntax-and-semantics
% Section: introduction

\documentclass[../../../include/open-logic-section]{subfiles}

\begin{document}

\olfileid{nml}{syn}{int}

\olsection{ভূমিকা}

মোডাল যুক্তিবিদ্যায় \emph{মোডাল বচন} এবং তাদের মধ্যকার অনুবর্তন সম্পর্ক
আলোচনা করা হয়। মোডাল বচনের কয়েকটি উদাহরণ:
\begin{enumerate}
\item $2+2=4$ হওয়া আবশ্যিক।
\item আগামীকাল বৃষ্টি হওয়া আবশ্যিকভাবে সম্ভব।
\item $!A$ হওয়া যদি আবশ্যিকভাবে সম্ভব হয়, তবে $!A$ হওয়া সম্ভব।
\end{enumerate}
সম্ভাব্যতা ও আবশ্যিকতাই একমাত্র মোডাল ধারণা নয়: অন্যান্য একস্থানিক সংযোজককেও
মোডাল সংযোজক বলা হয়। যেমন, ``$!A$ হওয়া উচিত'', ``ভবিষ্যতে $!A$ হবে'',
``ডানা জানে যে $!A$'' কিংবা ``ডানা বিশ্বাস করে যে $!A$''।

মোডাল যুক্তিবিদ্যার প্রথম আবির্ভাব অ্যারিস্টটলের \emph{De
  Interpretatione}-তে। তিনিই প্রথম লক্ষ করেন যে আবশ্যিকতা থেকে
সম্ভাব্যতা অনুসৃত হয়, কিন্তু বিপরীতটি হয় না; সম্ভাব্যতা ও আবশ্যিকতাকে
একে অপরের সাহায্যে সংজ্ঞায়িত করা যায়; $!A \land !B$ সম্ভবত সত্য হলে
$!A$ ও $!B$ প্রত্যেকটি সম্ভবত সত্য, কিন্তু বিপরীতটি হয় না; আর
$!A \to !B$ আবশ্যিক হলে, $!A$ আবশ্যিক হওয়া থেকে $!B$-এর
আবশ্যিকতাও অনুসৃত হয়।

মোডাল যুক্তিবিদ্যার প্রথম আধুনিক পদ্ধতি সি.~আই. লুইসের গবেষণায় দেখা যায়,
যার পরিণতি লুইস ও ল্যাংফোর্ডের \emph{Symbolic Logic} (1932)।
বস্তুগত শর্তবাচক দিয়ে অনুবর্তন প্রকাশে তাঁরা সন্তুষ্ট ছিলেন না:
``$!A$ থেকে $!B$ অনুসৃত হয়'' বোঝাতে $!A \lif !B$ যথেষ্ট নয়।
তাঁরা বরং ``$!A$ হলে $!B$, এটি আবশ্যিক'' হিসেবে অনুবর্তনকে
চিহ্নিত করার প্রস্তাব দেন, যার প্রতীক $!A \strictif !B$। বিভিন্ন
ধর্ম আলাদা করতে গিয়ে লুইস পাঁচটি পৃথক মোডাল ব্যবস্থা শনাক্ত করেন:
\Log{S1}, \ldots, \Log{S4}, \Log{S5}; শেষ দুটির ব্যবহার আজও রয়েছে।

লুইস ও ল্যাংফোর্ডের পদ্ধতি ছিল সম্পূর্ণ \emph{বাক্যগঠনগত}: তাঁরা
যুক্তিসঙ্গত স্বতঃসিদ্ধ ও নিয়ম বেছে নিয়ে সেগুলির সাহায্যে কী প্রমাণ করা
যায় তা পরীক্ষা করেন। দীর্ঘকাল অর্থতাত্ত্বিক পদ্ধতি অধরাই ছিল। পরে
রুডলফ কার্নাপ \emph{Meaning and Necessity} (1947)-এ \emph{অবস্থা-বিবরণ}
ধারণা ব্যবহার করে প্রথম চেষ্টা করেন। অবস্থা-বিবরণ বলতে পরমাণু বাক্যগুলির
একটি সংগ্রহ বোঝায়, যেগুলি ওই বিবরণে ``সত্য''। সত্যের সংজ্ঞা যেকোনো
বাক্য $!A$-তে প্রসারিত করে কার্নাপ $!A$-কে \emph{আবশ্যিকভাবে সত্য}
বলেন যদি তা সব অবস্থা-বিবরণে সত্য হয়। কিন্তু তাঁর পদ্ধতি \emph{পুনরাবৃত্ত}
মোডাল সংযোজক সামলাতে পারে না: ``সম্ভবত আবশ্যিকভাবে \ldots সম্ভবত $!A$''
রূপের বাক্য সব সময়ে সবচেয়ে ভিতরের মোডাল সংযোজকে নেমে আসে।

মোডাল অর্থতত্ত্বে বড় অগ্রগতি আসে সল ক্রিপকের ``A Completeness Theorem
in Modal Logic'' (JSL 1959) প্রবন্ধে। লাইবনিজের ধারণা অনুযায়ী কোনো
বক্তব্য সব সম্ভাব্য জগতে সত্য হলেই তা আবশ্যিকভাবে সত্য; ক্রিপকে এই ধারণা
থেকে শুরু করেন। তবে এতে কার্নাপের পদ্ধতির একই অসুবিধা থাকে: কোনো
জগৎ $w$-তে (বা অবস্থা-বিবরণ $s$-এ) বক্তব্যের সত্যতা আদৌ $w$-এর
ওপর নির্ভর করে না। তাই ক্রিপকে জগৎগুলির মধ্যে একটি \emph{অভিগম্যতা
সম্পর্ক} $R$ ধরেন। ``আবশ্যিকভাবে $!A$'' $w$-তে সত্য হবে তখনই,
যখন $w$ থেকে \emph{অভিগম্য} প্রতিটি জগৎ $w'$-তে $!A$ সত্য হবে।
এই পদ্ধতির কোনো রূপ অনুসরণকারী অর্থতত্ত্বকে ক্রিপকে অর্থতত্ত্ব বলে;
এটিই বিভিন্ন মোডাল যুক্তিবিদ্যার দ্রুত বিকাশ সম্ভব করেছে।

ক্রিপকে অর্থতত্ত্বে ব্যাখ্যা করলে মোডাল যুক্তিবিদ্যা \emph{সম্পর্কমূলক
গঠন} ``ভিতর থেকে'' কেমন দেখায় তা জানায়। সম্পর্কমূলক গঠন হলো
দ্বিপদী সম্পর্কসহ একটি সেট। যেমন, একটি শ্রেণির শিক্ষার্থীদের সেটকে
তাদের সামাজিক নিরাপত্তা নম্বর দিয়ে ক্রমান্বিত করলে একটি সম্পর্কমূলক গঠন
পাওয়া যায়। বাস্তবে বিভিন্ন ক্ষেত্রেই এমন গঠন দেখা যায়: জগতের
অবস্থাগুলির আপেক্ষিক সম্ভাব্যতা ছাড়াও কোনো ব্যক্তির জ্ঞানগত অবস্থাগুলি
জ্ঞানগত সম্ভাব্যতার সম্পর্কে যুক্ত হতে পারে; আবার কোনো গতিশীল ব্যবস্থার
অবস্থাগুলি অবস্থা-পরিবর্তনের সম্পর্কে যুক্ত হতে পারে। মোডাল যুক্তিবিদ্যা
দিয়ে এগুলির প্রত্যেকটির প্রতিরূপ তৈরি করা যায়। প্রথমটি দেয় সাধারণ
সত্যতা-সংক্রান্ত মোডাল যুক্তিবিদ্যা; অন্যগুলি দেয় জ্ঞানতাত্ত্বিক
যুক্তিবিদ্যা, গতিশীল যুক্তিবিদ্যা ইত্যাদি।

আমরা একটি বিশেষ দৃষ্টিকোণে মন দেব, মোডাল যুক্তিবিদ্যায় যার নাম
``অনুরূপতা তত্ত্ব''। ক্রিপকের প্রথম দিকের উল্লেখযোগ্য আবিষ্কারগুলির
একটি হলো: অভিগম্যতা সম্পর্ক~$R$-এর অনেক ধর্ম (যেমন পরিযায়িতা,
প্রতিসাম্য) যথাযথ ``মোডাল স্কিমা'' দিয়ে \emph{মোডাল ভাষার মধ্যেই}
চিহ্নিত করা যায়। যেমন, মোডাল যুক্তিবিদেরা বলেন, $R$-এর প্রতিবিম্ব ধর্ম
``$!A$ আবশ্যিক হলে $!A$'' স্কিমার সঙ্গে ``অনুরূপ''। আমরা মূলত
\Ax{D}, \Ax{T}, \Ax{B}, \Ax{4} ও~\Ax{5} স্কিমা মিলিয়ে পাওয়া
কয়েকটি ধ্রুপদি মোডাল ব্যবস্থার (যেমন \Log{S4} ও \Log{S5})
অনুরূপতা তত্ত্ব আলোচনা করব।

\end{document}
